\documentclass{article}

\usepackage{amsmath}
\usepackage{amssymb}
\usepackage{enumitem}
\usepackage{xcolor}

\title{Weekly Homework 03}
\author{Trevon Collins}
\date{}

\setcounter{secnumdepth}{0}
\begin{document}
\maketitle
\section{Textbook Problems}
\begin{enumerate}[label=\#\arabic*]
    \setcounter{enumi}{3}
    \item Suppose $x$, $y \in \mathbb{Z}$. If $x$ and $y$ are odd, then xy is odd.
    \begin{enumerate}[label=Case \arabic*:]
        \item $x > 0$ and $y > 0$
        \item $x < 0$ and $y > 0$
        \item $x > 0$ and $y < 0$
        \item $x < 0$ and $y < 0$
    \end{enumerate}
    Proof 1: Suppose $x=2a+1$, for all integers $a$, $x$ is odd. 
    W.L.O.G. we can say the same for $y$ for $y=2b+1$.
    \[
        \begin{aligned}
            &xy = (2a+1)(2b+1) = 4ab + 2a + 2b + 1 = 2(2ab + a + b) + 1\\
            &\implies 2m + 1 \text{, which is odd.}
        \end{aligned}
    \]

    \setcounter{enumi}{5}
    \item Suppose $a$, $b$, $c \in \mathbb{Z}$. If $a|b$ and $a|c$, then $a|(b + c)$.
    \begin{enumerate}[label=Case \arabic*:]
        \item $b = 2m$ and $c = 2n$
        \item $b = 2m$ and $c = 2n + 1$
        \item $b = 2m + 1$ and $c = 2n + 1$
    \end{enumerate}
    \begin{enumerate}[label=Proof \arabic*:]
        \item $b = a \times k$ and $c = a \times l$ should mean
        $b+c = a \times k + a \times l = a(k + l)$
        \[
            \begin{aligned}
                2m + 2n &= 2(\frac{ak}{2}) + 2(\frac{al}{2}) 
                = 2(\frac{ak}{2} + \frac{al}{2}) = ak + al = a(k + l)
            \end{aligned}
        \]
        \item W.L.O.G.
        \[
            \begin{aligned}
                &2m + 2n + 1 = 2(\frac{ak}{2}) + 2(\frac{al - 1}{2} + 1)
                = 2(\frac{ak}{2} + \frac{al - 1}{2}) + 1 \\
                &= ak + al - 1 + 1 = a(k + l)
            \end{aligned}
        \]
        \item W.L.O.G.
        \[
            \begin{aligned}
                &2m + 1 + 2n + 1 = 2(\frac{ak - 1}{2}) + 1 + 2(\frac{al - 1}{2}) + 1 \\
                &= 2(\frac{ak - 1}{2} + \frac{al - 1}{2}) + 2 = ak + al - 1 - 1 + 2 = a(k + l)
            \end{aligned}
        \]
    \end{enumerate}
    \item  Suppose $a, \ b \in \mathbb{Z}$. If $a|b$, then $a^2|b^2$.
    \begin{enumerate}[label=Case \arabic*:]
        \item If $a = 0$, then $b = 0$ and $a^2 | b^2$.
        \item If $a \neq 0$, then $b = ak$ for some integer $k$. Therfore $b^2 = a^2k^2$ and $a^2 | b^2$. W.L.O.G. applies here for all integers.
    \end{enumerate}

    \setcounter{enumi}{13}
    \item If $n \in \mathbb{Z}$, then $5n^2+3n+7$ is odd.
    \begin{enumerate}[label=Case \arabic*:]
        \item $n = 2k$ for some integer $k$.
        \item $n = 2k + 1$ for some integer $k$.
    \end{enumerate}
    \begin{enumerate}[label=Proof \arabic*:]
        \item
        \[
            \begin{aligned}
                &5n^2 + 3n + 7 = 5(2k)^2 + 3(2k) + 7 = 20k^2 + 6k + 7 \\
                &= 20k^2 + 6k + 6 + 1 = 2(10k^2 + 3k + 3) + 1\\
                &\implies 2m + 1 \text{, which is odd.}
            \end{aligned}
        \]
        \item 
        \[
            \begin{aligned}
                &5n^2 + 3n + 7 = 5(2k + 1)^2 + 3(2k + 1) + 7 \\
                &= 5(4k^2 + 4k + 1) + 6k + 3 + 7 = 20k^2 + 26k + 15 \\
                &= 20k^2 + 26k + 14 + 1 = 2(10k^2 + 13k + 7) + 1\\
                &\implies 2m + 1 \text{, which is odd.}
            \end{aligned}
        \]
    \end{enumerate}

    \setcounter{enumi}{15}
    \item  If two integers have the same parity, then their sum is even.
    \begin{enumerate}[label=Case \arabic*:]
        \item $k, \ l \in \mathbb{Z}$ and $k = 2m$ and $l = 2n$ for some integers $m$ and $n$.
        \item $k, \ l \in \mathbb{Z}$ and $k = 2m + 1$ and $l = 2n + 1$ for some integers $m$ and $n$.
    \end{enumerate}
    \begin{enumerate}[label=Proof \arabic*:]
        \item $k + l = 2m + 2n = 2(m + n)$, which is even.
        \item $k + l = 2m + 1 + 2n + 1 = 2(m + n + 1)$, which is even.
    \end{enumerate}
\end{enumerate}

\section{Additional Problems}
\begin{enumerate}
    \item Write problems in English sentences.
        \begin{enumerate}[label=\Alph*.]
            \item $\forall m \in \mathbb{Z}, \ \exists n \in \mathbb{Z}$ such that $m = 2n$. 
            \[
            \textcolor{blue}{\text{For integers $m$ there is an integer $n$ such that $m = 2n$.}}
            \]
            \item $\forall x \in \mathbb{R}, \ \exists y \in \mathbb{R}$ such that $x = 2y$.
            \[
            \textcolor{blue}{\text{For real numbers $x$ there is a real number $y$ such that $x = 2y$.}}
            \]
        \end{enumerate}
    \newpage
    \item An integer $x$ is trodd if $x = 3a + 1$ for some $a \in \mathbb{Z}$
        \begin{enumerate}[label=\Alph*.]
            \item
            \[
            \begin{array}{c|c}
                a & x = 3a + 1 \\ \hline
                0 & 1 \\
                1 & 4 \\
                2 & 7 \\
                3 & 10 \\
                4 & 13 
            \end{array}
            \]
            By this table, we can see that 2, 3, 5, 6, and 8 would not be trodd.
            \item If $x$ is a trodd integer, then $x^2$ is also a trodd integer.
            \[
                \begin{aligned}
                    \text{Proof: }&x^2 = (3a + 1)^2 = 9a^2 + 6a + 1 = 3(3a^2 + 2a) + 1 \\
                    &\implies x^2 \text{ is a trodd integer.}
                \end{aligned}
            \]
        \end{enumerate}
    \item An integer $x$ is four-six divisible if $4|x$ or $6|x$.
        \begin{enumerate}[label=\Alph*.]
            \item Prove: If $x$ is four-six divisible, then $7x + 3$ is odd.
            \begin{enumerate}[label=Case \arabic*:]
                \item If $4|x$, then $x = 4a$ for $a \in \mathbb{Z}$
                \item If $6|x$, then $x = 6b$ for $b \in \mathbb{Z}$
            \end{enumerate}
            \begin{enumerate}[label=Proof \arabic*:]
                \item $7x + 3 = 7(4a) + 3 = 28a + 3 = 2(14a) + 3$, which is odd.
                \item $7x + 3 = 7(6b) + 3 = 42b + 3 = 2(21b) + 3$, which is odd.
            \end{enumerate}
        \end{enumerate}
\end{enumerate}
\end{document}